\chapter[RELATÓRIO DE ATIVIDADES DESENVOLVIDAS]{Relatório de atividades desenvolvidas}

\begin{table}[htb]
    \centering
    \caption{\label{tab:tabularx} Resumo das atividades desenvolvidas durante o estágio.}

    \begin{tabularx}{\textwidth}{ l X }
        \hline
        \textbf{Semana} & \textbf{Atividades} \\ \hline
        22-24/abril & Acompanhamento de pedidos \\ 
        27-30/abril & Revisão de projetos arquitetônicos da prefeitura \\ 
        04-08/maio & Orçamentos \\ 
        11-14/maio & Modelagem 3D \\ 
        \hline
    \end{tabularx}
    
    \vspace{1em}
    \fonte{Elaborado pelo autor.}
\end{table}

**Descrição Minuciosa das Atividades:**
Durante os 57 dias de estágio, as atividades foram divididas em três frentes principais de atuação no escritório:

*   **Elaboração e Revisão de Projetos:**
    Fui responsável por auxiliar na elaboração, correção e revisão de projetos arquitetônicos e elétricos, lidando diretamente com as revisões solicitadas por clientes e exigências de aprovação da prefeitura (como o projeto do cliente José Henrique e projetos sob responsabilidade de Nelson Realce e Vicktor). As ferramentas utilizadas incluíram o software SketchUp para modelagem e a inserção de informações em plantas no AutoCAD (como na planta Padre Pio). Também atuei na parte de legalização e documentação técnica, elaborando Memoriais Descritivos, processos de Habite-se, Termos de Compromisso (incluindo termo de calçada) e organizando a dobragem de folhas A1 de projeto para entrega. 

*   **Orçamentação e Levantamento de Quantitativos:**
    Na área de custos, realizei o levantamento de quantitativos e a correção de orçamentos, incluindo a orçamentação de lajes e alterações de tipos de cobertura para cerâmica. Uma atividade de destaque foi o ajuste de orçamentos com a exclusão de serviços de gesso e a atualização sistemática dos valores utilizando a base de referência do SINAPI, atualizada para o mês 03/2026.
    
*   **Gestão de Suprimentos, Cotações e Rotinas Administrativas:**

    No suporte à gestão, realizei o controle de materiais preenchendo planilhas de acompanhamento de pedidos e de projetos, além de realizar *checklists*. Efetuei diversas cotações e pedidos de insumos diretos e indiretos para as obras e escritório, tais como: tintas, serviços de pintor, cascalho, pneus, filtros para enceradeira, polias de sondagem, telas e tanques para compactador/sapo. Fui encarregada também do controle financeiro básico, colocando notas fiscais nos pedidos, faturando pedidos de fornecedores (como Construlares) e fazendo a apropriação das NFs recebidas, vinculando-as corretamente ao pedido, etapa da obra e responsável.